\section{Convergence inputs and relative structure}
\label{sec:structure}

The auxiliary systems in the proof will be extensions and joinings of a
suspension of the original transformation.  They need not be ergodic.
We therefore formulate the structural facts in the nonergodic setting,
and distinguish the pointwise input for three functions from the norm
estimate available for any finite number of functions.

\subsection{Standard systems and the pointwise inputs}

Put $I=\{1,2,3,4\}$.  It suffices to prove
Theorem~\ref{thm:main} on a standard probability space.  Indeed, for a
given tuple of bounded functions, the map recording all their integer
translates takes values in a countable product of compact scalar discs.
The pushforward measure and the shift give a standard factor, the
functions descend to its coordinate functions, and mixing passes to
this factor.  Almost-sure convergence there pulls back to the original
space.  We henceforth write $(X,\mu,T)$ for this standard mixing system.
All factors and conditional expectations are understood modulo null
sets; probability spaces may be completed.

We use \cite[Theorem~1.1]{MixingCompanion} in the following
form.  For every fixed finite family $g_1,\ldots,g_s\in L^\infty(\mu)$,
\begin{equation}\label{eq:all-orders-input}
 \int_X\prod_{j=1}^s g_j(T^{n_j}x)\,\dd\mu(x)
 \longrightarrow \prod_{j=1}^s\int_Xg_j\,\dd\mu
 \quad\text{if}\quad \min_{j\ne l}|n_j-n_l|\longrightarrow\infty.
\end{equation}
The bounded-function formulation follows from the set formulation by
simple-function approximation.  All uses of this input involve fixed
functions; no rate of mixing is assumed.

The other companion input will be used on systems without mixing.
Its coefficient dependence is important, so we isolate its exact scope.

\begin{lemma}[Lower-length convergence]\label{lem:lower-length}
Let $(Y,m,R)$ be an invertible probability-preserving transformation on
a standard probability space.  If $\varnothing\ne J\subsetneq I$ and
$g_i\in L^\infty(m)$ for $i\in J$, then
\[
 \frac1N\sum_{n=1}^N\prod_{i\in J}g_i(R^{in}y)
\]
converges for $m$-almost every $y$.
\end{lemma}

The proof, including the changes of coefficients in the analytic
estimates of \cite{TripleCompanion,HilbertCompanion}, is given in
Appendix~\ref{sec:lower-length}.  Neither ergodicity nor mixing is an
assumption of Lemma~\ref{lem:lower-length}.

\subsection{Distal factors and relatively weakly mixing extensions}

Let $(Y,m,R)$ be a standard probability-preserving system, and let
$B$ be a factor.  We use the conditional-correlation characterization
of \emph{relative weak mixing}: for bounded $f,g$, with
$\E[f\mid B]=0$,
\begin{equation}\label{eq:rwm-correlations}
 \frac1H\sum_{h=1}^H
 \bigl\|\E[\overline f\,(g\circ R^h)\mid B]\bigr\|_2^2
 \longrightarrow0.
\end{equation}
For $F_1,\ldots,F_r\in L^2(m)$, the associated finitely generated
$L^\infty(B)$-module is
\[
 \left\{\sum_{j=1}^r a_j F_j:a_j\in L^\infty(B)\right\}.
\]
It is invariant if composition with $R$ and with $R^{-1}$ preserves
this space.  An extension is \emph{relatively compact} if its $L^2$
space is generated by these invariant modules.
The individual modules need not be closed in $L^2$; closure is taken
after their union is spanned.  This is the compact-module
characterization in its corrected form
\cite[Theorem~4.1]{Jamneshan}\cite[Theorem~2.5]{JamneshanCorrigendum}.
A
\emph{relatively distal} extension is a tower of relatively compact
extensions and inverse limits.  Distality without a specified base
means distality over the trivial factor.  These definitions include
nonergodic systems; in particular an identity action is compact.

\begin{proposition}[Relative structure]\label{prop:structure}
The following statements hold for standard probability-preserving
systems.
\begin{enumerate}
\item There is a maximal measure-distal factor $W_Y$.  It contains the
invariant factor, and $Y\to W_Y$ is relatively weakly mixing.  A distal
factor pulls back into the maximal distal factor of any extension.
Every measure-preserving automorphism commuting with $R$ preserves
$W_Y$.
\item Factors and joinings of distal systems are distal.  A distal
system is relatively distal over each of its factors.  For every
nonzero integer $q$, $R$ is distal if and only if $R^q$ is distal.
\item Relative weak mixing passes to nonzero integer powers and to
conditionally independent base change.  More generally, the product
of relatively weakly mixing fibre laws over any invariant joining of
their base laws is relatively weakly mixing over that joining.
\item A relatively weakly mixing extension and a relatively distal
extension of the same base are relatively disjoint: the only invariant
joining which identifies their base coordinates is the relatively
independent joining.
\item If $R$ is distal, then, for every $d\ge1$ and bounded
$g_1,\ldots,g_d$, the averages
\[
 \frac1N\sum_{n=1}^N\prod_{i=1}^d g_i(R^{in}y)
\]
converge almost surely.
\end{enumerate}
If $(S_Y^t)_{t\in\R}$ is a probability-preserving flow, the maximal
distal factor $W_Y$ of $S_Y^1$ is a flow factor.  Its action at every
nonzero rational time is distal.
\end{proposition}

\begin{proof}
The existence assertion and the relatively weakly mixing extension
are the Furstenberg--Zimmer structure theorem
\cite{Furstenberg,Zimmer}; its nonergodic probability-algebra formulation
is given in \cite[Theorems~1.1 and~1.2]{Jamneshan}.
We recall the consequences needed here, with particular attention to
changes of base and of time.

At a compact stage, an invariant finitely generated module remains such a
module after the base is enlarged.  One may use bounded generators:
for a vector of generators transforming by unitary matrices over the
base, truncate by its invariant pointwise length.  Products of bounded
generators lie in finitely generated tensor modules.  Consequently the join
of compact stages is compact over the joined base, and the same
assertion for distal towers follows by passing to closures at limit
stages.  This proves closure under joinings and the assertion about
enlarging the base to any factor of a distal system.  It also shows
that distal factors of a system lie in its maximal distal factor.
Applying this observation to an automorphism commuting with $R$, and
then to its inverse, shows that it preserves $W_Y$.  The invariant
factor is compact and therefore belongs to $W_Y$.

Compact modules for $R$ are compact modules for $R^q$, so distality
passes to powers.  Conversely, suppose $q>0$ and $R^q$ is distal.
Its canonical compact tower is preserved by $R$, which commutes with
$R^q$.  At each successor stage replace an $R^q$-invariant finitely
generated module by the sum of its first $q$ translates under $R$.
This is a finitely generated $R$-invariant module over the same base.
The resulting tower
proves distality for $R$.  Negative powers give the same conclusion
by inversion.

For relative weak mixing, the nonnegative summands in
\eqref{eq:rwm-correlations} show that restriction to multiples of a
fixed positive integer still has Ces\`aro limit zero.  Invariance and
complex conjugation handle negative powers.  To verify the assertion
about joined bases, let $Y_j\to B_j$ be the extensions and let $\sigma$
be an invariant joining of the base measures.  On the resulting space
give the fibres the law
\[
 \int\bigotimes_j m_j(\dd y_j\mid b_j)\,\dd\sigma((b_j)_j).
\]
For bounded tensor tests, their conditional correlation over the full
base has the form
\begin{equation}\label{eq:product-rwm-correlation}
 a(b)\,c(S_B^h b)
 \prod_j\E[\overline{f_j}\,(g_j\circ R_j^h)\mid B_j](b_j).
\end{equation}
If one $f_j$ is conditionally centred, the squared $L^2(\sigma)$ norm
is bounded by a constant times the squared norm of that factor in
$L^2(B_j)$.  Here only the $j$th marginal of $\sigma$ is used.
Its Ces\`aro mean tends to zero.  Subtracting the conditional
expectation of a tensor and telescoping shows that such centred
tensors span a dense subspace of the orthocomplement of the full
base.  Bounded approximation now proves the claim.  The same
calculation with an arbitrary extra base coordinate proves stability
under conditionally independent base change.

For relative disjointness we argue with conditional modules.  The
absolute weak-mixing disjointness principle goes back to
Furstenberg~\cite[Part~I, Section~7]{FurstenbergDisjointness}.
First take a compact extension of the base.  A joining defines a
base-linear, time-intertwining conditional
expectation operator from its $L^2$ space to the relatively weakly
mixing side.  It takes each finitely generated invariant module into
another such module.  Relative weak mixing says that the latter lies
in the base, so the joining is relatively independent.  Induct along
a distal tower: once independence has been proved at a stage, the
extension on the other side is the conditionally independent base
change of the original relatively weakly mixing extension.  The
preceding calculation allows the compact-stage argument to be
repeated.  Conditional expectations pass to the limit at limit
stages.  This proves relative disjointness without an ergodicity
assumption.  It also gives closure of distal systems under factors:
if $D$ is distal and $B$ is its factor, then $D$ is distal over the
pullback of $W_B$; relative disjointness of $D\to W_B$ and
$B\to W_B$ forces the latter extension to be trivial.

For the pointwise assertion, the ergodic case is
\cite[Theorem~C]{HuangShaoYe}, whose distality is measure-theoretic.
For a nonergodic distal system, start its compact tower with the
invariant factor and disintegrate over that factor.  The tower is
countable on a standard space: its genuinely increasing successor
stages give nonzero orthogonal differences in a separable $L^2$
space.  At each stage choose a countable generating algebra.  Approximate
each of its indicators by functions in the compact modules at a
successor stage, or in the preceding factors at a limit stage, with
summable squared $L^2$ errors.  Fubini makes these errors summable in
the conditional $L^2$ spaces on almost every ergodic component,
simultaneously for all the chosen approximations.  The restricted
generating algebras are dense for every component measure.  The
countably many identities expressing invariance of the chosen modules
also disintegrate.  Thus the successor stages remain compact and the
limit stages remain generated by their predecessors on a common
conull set of components.  Every intermediate factor has the same invariant
algebra, so each component tower consists of ergodic systems; its
compact extensions are isometric by the compact-extension structure
theorem.  Each component is thus an ergodic distal system.
The cited theorem applies to the given tuple on almost every
component, proving the assertion.

Finally, every flow time commutes with time one and therefore
preserves $W_Y$.  If $r=p/q\ne0$ is rational, the $q$th power of
its action on $W_Y$ is the distal action at integer time $p$.
The power assertion gives distality at time $r$.
\end{proof}

The following norm estimate explains why proper subaverages have a
particularly simple law over a distal factor.  Its proof uses relative
weak mixing alone.

\begin{lemma}[A relatively weakly mixing characteristic factor]
\label{lem:relative-norm}
Suppose $Y\to B$ is relatively weakly mixing for an invertible
transformation $R$.  If $a_1,\ldots,a_d$ are distinct nonzero integers
and $g_1,\ldots,g_d$ are bounded, then
\begin{equation}\label{eq:relative-norm}
 \left\|\frac1N\sum_{n=1}^N
 \left(\prod_{i=1}^d g_i\circ R^{a_i n}
 -\prod_{i=1}^d\E[g_i\mid B]\circ R^{a_i n}\right)\right\|_2
 \longrightarrow0.
\end{equation}
\end{lemma}

\begin{proof}
By telescoping and scaling, it suffices to consider functions bounded
by one, one of which, say $g_s$, has conditional mean zero.  We induct
on $d$, with the empty product as the initial case.  Set
$v_n=\prod_i g_i\circ R^{a_i n}$ and, for a fixed $h\ge1$, put
$F_{i,h}=\overline{g_i}(g_i\circ R^{a_i h})$.
Invariance gives
\[
 \langle v_n,v_{n+h}\rangle
 =\int F_{s,h}\prod_{i\ne s}
       F_{i,h}\circ R^{(a_i-a_s)n}\,\dd m.
\]
The other coefficients are distinct and nonzero.  The induction
hypothesis replaces their factors, in the average over $n$, by
their conditional expectations over $B$, with an error tending to
zero.  Those remaining factors are $B$-measurable and bounded by
one, so
\[
 \limsup_{N\to\infty}
 \left|\frac1N\sum_{n=1}^N\langle v_n,v_{n+h}\rangle\right|
 \le \|\E[F_{s,h}\mid B]\|_1.
\]
For $d=1$ the same inequality holds directly.  The Hilbert-space
van der Corput inequality bounds the squared norm of the averages
of $v_n$ by an absolute constant times
\[
 \frac1H+\frac1H\sum_{h=1}^H\|\E[F_{s,h}\mid B]\|_1
\]
in the limit superior as $N\to\infty$.  The second term tends to
zero by \eqref{eq:rwm-correlations} for $R^{a_s}$ and
Cauchy--Schwarz.  Letting $H\to\infty$ proves the induction.
\end{proof}

In particular, weak mixing of $T$ gives $L^2$ convergence to the
product of the means for all finite distinct nonzero linear powers
on $X$.  In a flow, with $B$ preserved by every time, the analogue of
\eqref{eq:relative-norm} for deterministic continuous-time averages
follows by writing $t=n+s$, $0\le s<1$.  Apply the lemma for each
fixed $s$ to $g_i\circ S_Y^{a_i s}$, and integrate the resulting
$L^2$ estimate in $s$.  Conditional expectation commutes with these
time shifts.  Bounded convergence and the final interval of length
less than one complete the passage to continuous time.
